\documentclass[11pt,a4paper]{article}
\usepackage{amsmath,amssymb,amsthm,physics,geometry}
\usepackage{mathtools}
\usepackage{hyperref}
\geometry{margin=1in}

\title{\textbf{Inverse–Square Potential on the Half–Line\\
and Conformal Quantum Mechanics}}
\author{}
\date{}

\begin{document}
\maketitle

\begin{abstract}
We connect the one–dimensional half–line inverse–square
Schrödinger Hamiltonian with the generators of conformal quantum mechanics.
The Hamiltonian, dilation, and special conformal operators
form the Lie algebra $\mathfrak{so}(2,1)$.
We analyze the representation parameter,
its relation to the coupling constant,
and show how the hard–wall boundary condition arises
as a singular strong–coupling limit within the conformal algebra.
\end{abstract}

\section{Hamiltonian on the Half–Line}

Consider
\begin{equation}
H
=
\frac{p^2}{2m}
+
\frac{g}{x^2},
\qquad x>0,
\qquad p=-i\hbar\frac{d}{dx}.
\end{equation}

Introduce
\begin{equation}
\gamma=\frac{2mg}{\hbar^2},
\qquad
\nu=\sqrt{\frac{1}{4}+\gamma}.
\end{equation}

The stationary equation becomes
\begin{equation}
\psi''(x)+\left(k^2-\frac{\gamma}{x^2}\right)\psi(x)=0.
\end{equation}

For $g>0$, the Hamiltonian is essentially self–adjoint,
and the physical solution is
\begin{equation}
\psi(x)=\sqrt{x}J_\nu(kx).
\end{equation}

\section{Conformal Generators}

Define the three operators

\begin{align}
H &= \frac{p^2}{2m}+\frac{g}{x^2}, \\
D &= \frac{1}{2}(xp+px), \\
K &= \frac{m x^2}{2}.
\end{align}

They satisfy the commutation relations

\begin{align}
[ D,H ] &= i\hbar H, \\
[ D,K ] &= -i\hbar K, \\
[ H,K ] &= 2i\hbar D.
\end{align}

This is the Lie algebra $\mathfrak{so}(2,1)$.

\section{Casimir Operator}

The quadratic Casimir is
\begin{equation}
C = HK + KH - D^2.
\end{equation}

A direct calculation yields
\begin{equation}
C
=
\frac{\hbar^2}{4}
\left(
\nu^2-\frac{1}{4}
\right).
\end{equation}

Thus the representation is characterized by
\begin{equation}
\nu=\sqrt{\frac{1}{4}+\frac{2mg}{\hbar^2}}.
\end{equation}

Hence:

\begin{center}
The coupling constant $g$ determines the
$\mathfrak{so}(2,1)$ representation weight.
\end{center}

\section{Dilation Properties}

The dilation generator acts as
\begin{equation}
D=-i\hbar\left(x\frac{d}{dx}+\frac{1}{2}\right).
\end{equation}

The regular solution behaves near $x=0$ as
\begin{equation}
\psi(x)\sim x^{\nu+\frac{1}{2}}.
\end{equation}

Thus it is an eigenfunction of $D$ asymptotically:
\begin{equation}
D\,x^{\nu+\frac{1}{2}}
=
-i\hbar\left(\nu+1\right)x^{\nu+\frac{1}{2}}.
\end{equation}

The exponent $\nu+\frac{1}{2}$ is the conformal scaling dimension.

\section{Compact Generator (DFF Construction)}

Define the compact Hamiltonian
\begin{equation}
H_\omega = H+\omega^2 K.
\end{equation}

Explicitly,
\begin{equation}
H_\omega
=
\frac{p^2}{2m}
+
\frac{g}{x^2}
+
\frac{1}{2} m\omega^2 x^2.
\end{equation}

This has discrete spectrum
\begin{equation}
E_n
=
2\hbar\omega
\left(
n+\frac{1}{2}+\frac{\nu}{2}
\right),
\quad n=0,1,2,\dots
\end{equation}

The operators
\begin{equation}
J_0=\frac{1}{2\omega}H_\omega,
\qquad
J_\pm
=
\frac{1}{2}\left(\frac{K}{\omega}-\omega H\right)
\mp iD
\end{equation}
satisfy
\begin{align}
[J_0,J_\pm]&=\pm J_\pm,\\
[J_+,J_-]&=-2J_0.
\end{align}

Thus the spectrum is generated algebraically.

\section{Hard Wall as Strong–Coupling Limit}

As $g\to+\infty$:
\begin{equation}
\nu\sim \sqrt{\frac{2mg}{\hbar^2}}\to\infty.
\end{equation}

Then near the origin
\begin{equation}
\psi(x)\sim x^{\nu+\frac{1}{2}}.
\end{equation}

For any fixed $x>0$:
\begin{equation}
\lim_{g\to\infty} x^{\nu+\frac{1}{2}}=0.
\end{equation}

In representation terms:
\begin{itemize}
\item The Casimir $C\to\infty$.
\item The conformal weight diverges.
\item The scaling dimension becomes infinite.
\end{itemize}

This forces
\begin{equation}
\psi(0)=0,
\end{equation}
which is precisely the Dirichlet hard–wall boundary condition.

Thus:

\begin{center}
The hard wall is the infinite–weight limit
of a conformal $\mathfrak{so}(2,1)$ representation.
\end{center}

\section{Structural Picture}

\begin{itemize}
\item $H$ generates time translations.
\item $D$ generates scale transformations.
\item $K$ generates special conformal transformations.
\item $g$ fixes the representation via the Casimir.
\item $g\to\infty$ collapses the representation into
a boundary constraint.
\end{itemize}

Therefore the half–line inverse–square system
is the minimal conformal quantum mechanics model,
with the hard wall emerging as a singular representation limit.

\section{Conclusion}

The inverse–square Hamiltonian on the half–line:

\begin{enumerate}
\item Realizes $\mathfrak{so}(2,1)$ exactly.
\item Has coupling–dependent conformal weight $\nu$.
\item Becomes a hard wall when $\nu\to\infty$.
\item Provides the primitive building block of
conformal quantum mechanics.
\end{enumerate}

\end{document}
